\label{chap:methodology}
This chapter describes experimental designs and implementation details of this work. What libraries, hardware, datasets and procediments were used to achieve the results in chapter \ref{chap:results}. The goal of this study was to train and compare different \gls{DL} architectures to binary classificate images samples from \gls{CL} in order to detect \gls{CRC}.

\section{Datasets}
\label{sec_datasets}
For this survey, three different datasets were used: one for train, validation and test and the others to test only in the final stage of the 

\subsection{NCT-CRC-HE-100K}
\label{subsec_nct100k}
The main \gls{DS} of this study was 100,000 histological images of human colorectal cancer and healthy tissue, by Kather et al. \cite{kather2018crc100k}, samples were collected from the NCT Biobank (National Center for Tumor Diseases, Heidelberg, Germany) and the UMM pathology archive (University Medical Center Mannheim, Mannheim, Germany). All images were 224 x 224 pixels at 0.5 MMP and classified as one of the nine classes above:
\begin{itemize}
    \item Adipose (ADI)
    \item ackground (BACK)
    \item debris (DEB)
    \item lymphocytes (LYM)
    \item mucus (MUC)
    \item smooth muscle (MUS)
    \item normal colon mucosa (NORM)
    \item cancer-associated stroma (STR)
    \item colorectal adenocarcinoma epithelium (TUM)
\end{itemize}

They already divided the \gls{DS} in three differents sets:
\begin{itemize}
    \item NCT-CRC-HE-100K: contains one hundred thousand images color-normalized using Macenko's method \cite{5193250}.
    \item CRC-VAL-HE-7K: contains 7180 image from 50 different patients with no overlap with the ones in "NCT-CRC-HE-100K". The classes of this dataset are:
    \item Adipose (ADI)
    \item Background (BACK)
    \item debris (DEB)
    \item lymphocytes (LYM)
    \item mucus (MUC)
    \item smooth muscle (MUS)
    \item normal colon mucosa (NORM)
    \item cancer-associated stroma (STR)
    \item colorectal adenocarcinoma epithelium (TUM)
    \item NCT-CRC-HE-100K-NONORM: created from the same raw data as "NCT-CRC-HE-100K", however with no color normalization applied.
\end{itemize}


\subsection{Collection of textures in colorectal cancer histology}
\label{subsec_crc5000}
Collection of textures in histological images of human colorectal cancer by Kather et al. \cite{kather_2016_53169}, with two main folders:
\begin{itemize}
    \item Kather\_texture\_2016\_image\_tiles\_5000: Contains 5000 histological images of 150 x 150 px each.
    \begin{itemize}
        \item TUMOR
        \item STROMA
        \item COMPLEX
        \item LYMPHO
        \item DEBRIS
        \item MUCOSA
        \item ADIPOSE
        \item EMPTY
    \end{itemize}
    \item Kather\_texture\_2016\_larger\_images\_10: Contains 10 larger histological images of 5000 x 5000 px each, containing more than one tissue type.
\end{itemize}
 
\subsection{Lung and Colon Cancer Histopathological Image
Dataset (LC25000)}
\label{subsec_lc25000}
Composed of twenty five thousand images of colon and lung tissues, the LC25000 were created from originallt seven hundred fifty total images of lung tissue (250 benign lung tissue, 250 lung adenocarcinomas, and 250 lung squamous cell carcinomas) and 500 total images of colon tissue (250 benign colon tissue and 250 colon adenocarcinomas). All images were cropped to square sizes of 768 x 768 pixels from original 1024 x 768 pixels and augmented using the Augmentor software package from python programming language, by applying: left and right rotations (up to 25 degrees,1.0 probability) and by horizontal and vertical flips (0.5 probability).
The classes in the datasets are:
\begin{itemize}
    \item colon\_aca: colon adenocarcinomas
    \item colon\_n: benign colonic tissues
    \item lung\_aca: lung adenocarcinomas
    \item lung\_scc: lung squamous cell carcinomas
    \item lung\_n: benign lung tissues
\end{itemize}
All images were equally distributed between classes above, with five thousand for each one.

